\documentclass[a4paper,12pt]{article}
\usepackage[utf8]{inputenc}
\usepackage[T1]{fontenc}
\usepackage[french]{babel}
\usepackage{geometry}
\usepackage{booktabs}
\usepackage{array}
\usepackage{tabularx}
\usepackage{xcolor}
\usepackage{titlesec}
\usepackage{enumitem}
\usepackage{fancyhdr}
\usepackage{hyperref}

\geometry{margin=2.5cm}

\definecolor{primary}{RGB}{30, 90, 160}
\definecolor{secondary}{RGB}{230, 240, 255}
\definecolor{accent}{RGB}{200, 50, 50}

\titleformat{\section}{\large\bfseries\color{primary}}{}{0em}{}[\titlerule]
\titleformat{\subsection}{\normalsize\bfseries\color{primary}}{}{0em}{}

\pagestyle{fancy}
\fancyhf{}
\rhead{\textcolor{primary}{AP2 — Contexte M2L}}
\lhead{\textcolor{gray}{BTS SIO — SISR}}
\cfoot{\thepage}

\hypersetup{colorlinks=true, linkcolor=primary, urlcolor=primary}

\begin{document}

% ─── Titre ───────────────────────────────────────────────────────────────────
\begin{center}
    {\LARGE\bfseries\color{primary} AP2 — Contexte Maison des Ligues de Lorraine}\\[0.4em]
    {\large\color{gray} Récapitulatif SISR — PPE2}\\[0.2em]
    {\small BTS SIO | Réseau CERTA | Version 1.1}
\end{center}
\vspace{1em}

% ─── 1. Présentation M2L ─────────────────────────────────────────────────────
\section{Présentation de la M2L}

La \textbf{Maison des Ligues de Lorraine (M2L)} est un établissement du Conseil Régional de Lorraine,
cogéré avec le \textbf{CROSL} (Comité Régional Olympique et Sportif de Lorraine).
Sa mission est de fournir des espaces et des services aux ligues sportives régionales et structures hébergées.

\subsection{Structure physique}
\begin{itemize}[noitemsep]
    \item 4 bâtiments : A et C (4 étages), B et D (plain-pied)
    \item 80 bureaux pour 24 structures hébergées (ligues sportives, comités départementaux, CROSL)
    \item Espaces mutualisés : amphithéâtre 200 places, 8 salles de réunion (12–50 places),
          salle multimédia 24 postes, salle de convivialité, salle de reprographie
\end{itemize}

\subsection{Réseau informatique}

\begin{tabularx}{\textwidth}{>{\bfseries}lX}
    \toprule
    Zone & Équipement \\
    \midrule
    Bâtiment A & 4 armoires secondaires (1/étage), switches 26 ports (24×100M + 2×Giga), câblage Cat5 \\
    Bâtiment C & 1 armoire centrale (2\ieme\ étage), 4 switches 48 ports tout-Giga, câblage Cat6, liaison fibre vers armoire B \\
    Armoire principale (Bat. B) & 2 switches 24 ports + routeur + armoire serveurs (salle climatisée) \\
    \bottomrule
\end{tabularx}

\subsection{Services proposés}
\begin{itemize}[noitemsep]
    \item Internet, Wifi (bornes en quinconce dans tous les étages + RDC), Téléphonie IP
    \item Serveur FTP documentaire (utilisé à seulement 10\% de sa capacité)
    \item Intranet/Extranet : réservation salles (MRBS), annuaire des ligues, presse sportive régionale
    \item Réseau d'affichage numérique (3 écrans 42\textquotedbl\ dans les espaces de circulation)
    \item Système de gestion des configurations des postes ligues
    \item Affranchissement et reprographie numérique (facturés par le CROSL)
\end{itemize}

% ─── 2. Projets PPE2 ─────────────────────────────────────────────────────────
\section{Projets PPE2 disponibles (SISR)}

\begin{tabularx}{\textwidth}{>{\bfseries}l>{\itshape}lX}
    \toprule
    Projet & Durée & Résumé \\
    \midrule
    ASSIZ    & 12h & Couverture Wifi d'une manifestation sous chapiteau — VLANs, DHCP, NAT, filtrage \\
    DYLAN    & 8h  & Insertion dynamique d'un portable dans le VLAN de sa ligue par adresse MAC \\
    VELANNE  & —   & Segmentation complète du réseau M2L en VLANs (peu cadré) \\
    DEPLOG   & —   & Déploiement de logiciels à distance via gestionnaire de configs \\
    STATIP   & —   & DHCP statique : IP fixe assignée par adresse MAC \\
    SUPERVIZ & —   & Supervision des éléments actifs et gestion d'incidents \\
    \bottomrule
\end{tabularx}

\vspace{0.5em}
\colorbox{secondary}{\parbox{\dimexpr\textwidth-2\fboxsep}{%
    \textbf{\color{primary}Recommandation SISR :} choisir \textbf{ASSIZ + DYLAN}.
    Les deux projets couvrent l'essentiel du référentiel SISR2 sans se répéter,
    sont bien documentés avec des livrables clairs, et sont valorisants à l'oral.
}}

% ─── 3. Projet ASSIZ ─────────────────────────────────────────────────────────
\section{Projet ASSIZ — Couverture Wifi des Assises de l'Escrime}

\subsection{Contexte}
La ligue d'Escrime organise une manifestation sous chapiteaux (2500 visiteurs attendus) à quelques centaines
de mètres de la M2L. Le service informatique doit prototyper l'infrastructure réseau nomade.

\subsection{Architecture à mettre en place}

\begin{itemize}[noitemsep]
    \item \textbf{2 SSIDs} sur chaque borne Wifi :
    \begin{itemize}[noitemsep]
        \item \texttt{"public"} — diffusé, WPA2 (clé sur le badge visiteur), isolation entre postes
        \item \texttt{"orga"} — non diffusé, WPA2 (clé connue des organisateurs uniquement)
    \end{itemize}
    \item \textbf{2 VLANs} :
    \begin{itemize}[noitemsep]
        \item VLAN 2 → réseau orga
        \item VLAN 3 → réseau public
    \end{itemize}
    \item \textbf{DHCP} sur le routeur MDL :
    \begin{itemize}[noitemsep]
        \item Plage \texttt{10.10.10.0/24} pour le SSID orga
        \item Plage \texttt{172.31.1.0/24} pour le SSID public
    \end{itemize}
    \item \textbf{NAT} sur le routeur MDL (le routeur France Télécom ne peut pas être configuré)
    \item \textbf{Filtrage} : les postes du réseau public ne doivent pas accéder aux ressources du réseau orga
          (imprimante réseau, traceur à banderoles, afficheur géant)
\end{itemize}

\subsection{Éléments techniques clés}
\begin{itemize}[noitemsep]
    \item Les bornes Wifi et le switch s'administrent dans le \textbf{VLAN 1} par défaut
    \item Le switch et les bornes doivent supporter la gestion de VLANs associés à des SSIDs
    \item On peut choisir un routeur + switch séparés, ou un seul routeur/switch combiné
\end{itemize}

\subsection{Livrables attendus}
\begin{enumerate}[noitemsep]
    \item Démonstration du prototype fonctionnel
    \item Notice technique : configuration borne Wifi + switch + routeur
    \item Fichier de configuration du routeur sauvegardé sur serveur TFTP
\end{enumerate}

\subsection{Compétences SISR2 mobilisées}
\begin{itemize}[noitemsep]
    \item Configurer une maquette ou un prototype pour valider une solution
    \item Configurer les éléments d'interconnexion permettant la séparation des flux
    \item Configurer les éléments d'interconnexion permettant d'établir des périmètres de sécurité
    \item Configurer les éléments d'interconnexion permettant d'assurer la connexion avec les réseaux externes
    \item Valider et documenter une solution — Justifier le choix d'une solution technique
\end{itemize}

% ─── 4. Projet DYLAN ─────────────────────────────────────────────────────────
\section{Projet DYLAN — Insertion dynamique dans le VLAN par adresse MAC}

\subsection{Contexte}
Depuis la mise en place des VLANs (projet VELANNE), les personnels des ligues ne peuvent plus accéder
à leurs données depuis la salle de réunion commune de l'étage. Il faut qu'un PC portable se place
\textbf{automatiquement dans le bon VLAN} selon l'adresse MAC de la machine connectée.

\subsection{Architecture à mettre en place}

\begin{itemize}[noitemsep]
    \item Sur un switch \textbf{Cisco SF-300} (bâtiments A et B) :
    \begin{enumerate}[noitemsep]
        \item Créer les VLANs 121 (volley), 122 (escrime), 123 (badminton) — numérotation : bâtiment + étage + ligue
        \item Mettre les ports bureaux en mode \textbf{ACCESS} dans leur VLAN respectif
        \item Mettre le port de la prise de la salle de réunion en mode \textbf{GÉNÉRAL} (partagé entre les 3 VLANs)
        \item Déclarer des \textbf{groupes d'adresses MAC} (un groupe par ligue) via \texttt{VLAN Management > MAC-Based Groups}
        \item Mapper chaque groupe MAC vers son VLAN via \texttt{Mapping Group to VLAN}
    \end{enumerate}
\end{itemize}

\subsection{Principe de fonctionnement}
Quand un portable de la ligue de volley (adresse MAC dans le groupe 121) se branche sur la prise
de la salle de réunion, le switch l'affecte \textbf{automatiquement} au VLAN 121.
Il peut alors communiquer avec les postes fixes de la ligue de volley.

\subsection{Numérotation des VLANs M2L}
Format : \textbf{[Bâtiment][Étage][Numéro ligue]} — ex : VLAN 121 = bâtiment A, étage 2, ligue 1 (volley)

\subsection{Livrables attendus}
\begin{enumerate}[noitemsep]
    \item Démonstration du prototype (affectation MAC → VLAN)
    \item Notice technique : implémentation VLAN par adresse MAC sur Cisco SF-300
    \item Document avantages/inconvénients de la solution
    \item Support de formation : partage de dossiers Windows entre postes de deux machines du même VLAN
\end{enumerate}

\subsection{Compétences SISR2 mobilisées}
\begin{itemize}[noitemsep]
    \item Configurer une maquette ou un prototype pour valider une solution
    \item Configurer les éléments d'interconnexion permettant la séparation des flux
    \item Configurer les éléments d'interconnexion permettant d'établir des périmètres de sécurité
    \item Valider et documenter une solution — Justifier le choix d'une solution technique
\end{itemize}

% ─── 5. Pourquoi ce combo ────────────────────────────────────────────────────
\section{Pourquoi le combo ASSIZ + DYLAN est fort}

\begin{tabularx}{\textwidth}{lXX}
    \toprule
    & \textbf{ASSIZ} & \textbf{DYLAN} \\
    \midrule
    Thème central    & Wifi + sécurité réseau nomade & VLANs dynamiques par MAC \\
    Compétence clé   & NAT, filtrage, DHCP, VLANs    & VLAN-per-MAC sur Cisco \\
    Matériel         & Borne Wifi + switch + routeur  & Switch Cisco SF-300 \\
    Durée estimée    & 12h                            & 8h \\
    Originalité      & Fort (infrastructure nomade)   & Très fort (concept peu courant) \\
    \bottomrule
\end{tabularx}

\vspace{0.8em}
Les deux projets se \textbf{complètent sans se répéter} :
ASSIZ couvre la couche réseau complète (DHCP, NAT, filtrage, Wifi),
DYLAN approfondit les VLANs avec une mécanique plus avancée (affectation dynamique par MAC).
Ensemble ils couvrent presque tout le référentiel \textbf{SISR2}.

À l'oral, la progression est naturelle et cohérente :
\textit{« J'ai d'abord conçu une infrastructure Wifi sécurisée pour une manifestation nomade,
puis j'ai approfondi les VLANs en mettant en place une affectation dynamique par adresse MAC. »}

% ─── 6. Rappel modules SISR ──────────────────────────────────────────────────
\section{Rappel — Modules SISR exploitables sur M2L}

\begin{tabularx}{\textwidth}{>{\bfseries}lX}
    \toprule
    Module & Activité M2L correspondante \\
    \midrule
    SI1 & Intégration PC de ligue, antivirus, client FTP, image de restauration \\
    SI2 & Création VLAN ligue, modification affectation bureaux, analyse appel d'offres switches \\
    SI3 & Consultation BDD réservation salles, extraction/modification données \\
    SI5 & DHCP + relais, FTP sécurisé, sauvegarde automatique \\
    SISR1 & Assistance à distance, restauration poste de ligue \\
    SISR2 & VLANs, VPN, Wifi double réseau, tolérance de panne, VLAN dynamique \\
    SISR3 & Proxy avec priorisation, logs d'activité, haute disponibilité LDAP/Web \\
    SISR4 & Virtualisation salle formation, contrôleur de domaine secondaire \\
    SISR5 & Supervision éléments actifs M2L (switches, serveurs) \\
    SI7 & Logiciel de gestion d'incidents et des configurations \\
    \bottomrule
\end{tabularx}

\end{document}
